% GENERATED FILE. Edit canonical lessons, then run npm run generate:books.
% canonical-source-hash: fnv1a64:0f3fa18a251c4b6d
% canonical-lessons: BN-C541-protikkha, BN-C541-dirgho, BN-C541-songrokkhit, BN-C541-token, BN-C541-porjay

\chapter{Waiting Your Turn}
\label{ch:bn-waiting-your-turn}
\hlchaptermodality{541}

\begin{chapteropening}
\textbf{By the end of this chapter:} \emph{I can talk about waiting, rows, tokens and the unexpected.}

Complete the last lesson of chapter 541: I can talk about waiting, rows, tokens and the unexpected.
\end{chapteropening}

\section[\texorpdfstring{\bn{প্রতীক্ষা} (prôtīkkhā)}{প্রতীক্ষা (prôtīkkhā)}]{\bn{প্রতীক্ষা} (prôtīkkhā) — waiting; expectation}

\label{lesson:BN-C541-protikkha}



\begin{quote}
Before the new one: say the Bengali for armour, then the Bengali for to finish off.
\end{quote}

\subsection*{You'll want to know: \bn{প্রতীক্ষা}}
\textbf{\bn{প্রতীক্ষা}} — \emph{prôtīkkhā} — \textquotedblleft{}waiting; expectation\textquotedblright{}.

\begin{quote}
\bn{এক} \bn{বছরের} \bn{প্রতীক্ষার} \bn{পরে} \bn{ছেলেটা} \bn{বাড়ি} \bn{ফিরল।}
\end{quote}

\begin{quote}
\emph{Ek bôchhôrer prôtīkkhār pôre chheleṭā bāṛi phirlô.}
\end{quote}

\begin{quote}
After a year of waiting the boy came home.
\end{quote}

\begin{cousinweb}
\textbf{\bn{প্রতীক্ষা}} (\emph{prôtīkkhā}) is a tatsama, Sanskrit \emph{pratīkṣā}, \textquotedblleft{}expectation\textquotedblright{} (Monier-Williams): the longer, hopeful kind of waiting.
\end{cousinweb}

\begin{grammarlens}[title={What you've built}]
One more word for saying what happens if.
\end{grammarlens}

\subsection*{Your turn}
\begin{itemize}
\raggedright
  \item \emph{Say it:} \emph{prôtīkkhā}
  \item \emph{Say it:} \emph{prôtīkkhā}, once more
  \item \emph{Say it:} \emph{Ek bôchhôrer prôtīkkhār pôre chheleṭā bāṛi phirlô.} — After a year of waiting the boy came home.
  \item \emph{Recall:} say the Bengali for armour, then the Bengali for to finish off, then say \emph{prôtīkkhā} again
\end{itemize}

\begin{tcolorbox}[breakable,colback=teal!4,colframe=teal!35!black,title={Before you move on}]
What does \bn{প্রতীক্ষা} mean? (\textquotedblleft{}Waiting; expectation\textquotedblright{}.) Say it once more.
\end{tcolorbox}
\section[\texorpdfstring{\bn{দীর্ঘ} (dīrghô)}{দীর্ঘ (dīrghô)}]{\bn{দীর্ঘ} (dīrghô) — long-lasting; long (in time or length)}

\label{lesson:BN-C541-dirgho}



\begin{quote}
Before the new one: say the Bengali for to finish off, then the Bengali for expectation.
\end{quote}

\subsection*{You'll want to know: \bn{দীর্ঘ}}
\textbf{\bn{দীর্ঘ}} — \emph{dīrghô} — \textquotedblleft{}long-lasting; long (in time or length)\textquotedblright{}.

\begin{quote}
\bn{দীর্ঘ} \bn{লাইন} \bn{থাকলে} \bn{আমরা} \bn{কাল} \bn{আবার} \bn{আসব।}
\end{quote}

\begin{quote}
\emph{Dīrghô lāin thākle āmrā kāl ābār āsbô.}
\end{quote}

\begin{quote}
If there's a long queue, we'll come back tomorrow.
\end{quote}

\begin{cousinweb}
\textbf{\bn{দীর্ঘ}} (\emph{dīrghô}) is a tatsama, Sanskrit \emph{dīrgha}, \textquotedblleft{}long\textquotedblright{} (Monier-Williams), the formal partner of everyday \bn{লম্বা}.
\end{cousinweb}

\begin{grammarlens}[title={What you've built}]
One more word for saying what happens if.
\end{grammarlens}

\subsection*{Your turn}
\begin{itemize}
\raggedright
  \item \emph{Say it:} \emph{dīrghô}
  \item \emph{Say it:} \emph{dīrghô}, once more
  \item \emph{Say it:} \emph{Dīrghô lāin thākle āmrā kāl ābār āsbô.} — If there's a long queue, we'll come back tomorrow.
  \item \emph{Recall:} say the Bengali for to finish off, then the Bengali for expectation, then say \emph{dīrghô} again
\end{itemize}

\begin{tcolorbox}[breakable,colback=teal!4,colframe=teal!35!black,title={Before you move on}]
What does \bn{দীর্ঘ} mean? (\textquotedblleft{}Long-lasting; long (in time or length)\textquotedblright{}.) Say it once more.
\end{tcolorbox}
\section[\texorpdfstring{\bn{সংরক্ষিত} (sôngrôkkhit)}{সংরক্ষিত (sôngrôkkhit)}]{\bn{সংরক্ষিত} (sôngrôkkhit) — reserved; kept aside}

\label{lesson:BN-C541-songrokkhit}



\begin{quote}
Before the new one: say the Bengali for expectation, then the Bengali for long-lasting.
\end{quote}

\subsection*{You'll want to know: \bn{সংরক্ষিত}}
\textbf{\bn{সংরক্ষিত}} — \emph{sôngrôkkhit} — \textquotedblleft{}reserved; kept aside\textquotedblright{}.

\begin{quote}
\bn{ট্রেনের} \bn{এই} \bn{আসনগুলো} \bn{মহিলাদের} \bn{জন্য} \bn{সংরক্ষিত।}
\end{quote}

\begin{quote}
\emph{Ṭrener ei āsôngulo môhilāder jônyô sôngrôkkhit.}
\end{quote}

\begin{quote}
These seats on the train are reserved for women.
\end{quote}

\begin{cousinweb}
\textbf{\bn{সংরক্ষিত}} (\emph{sôngrôkkhit}) is a tatsama, Sanskrit \emph{saṃrakṣita}, \textquotedblleft{}protected, kept safe\textquotedblright{} (Monier-Williams): kept aside for someone.
\end{cousinweb}

\begin{grammarlens}[title={What you've built}]
One more word for saying what happens if.
\end{grammarlens}

\subsection*{Your turn}
\begin{itemize}
\raggedright
  \item \emph{Say it:} \emph{sôngrôkkhit}
  \item \emph{Say it:} \emph{sôngrôkkhit}, once more
  \item \emph{Say it:} \emph{Ṭrener ei āsôngulo môhilāder jônyô sôngrôkkhit.} — These seats on the train are reserved for women.
  \item \emph{Recall:} say the Bengali for expectation, then the Bengali for long-lasting, then say \emph{sôngrôkkhit} again
\end{itemize}

\begin{tcolorbox}[breakable,colback=teal!4,colframe=teal!35!black,title={Before you move on}]
What does \bn{সংরক্ষিত} mean? (\textquotedblleft{}Reserved; kept aside\textquotedblright{}.) Say it once more.
\end{tcolorbox}
\section[\texorpdfstring{\bn{টোকেন} (ṭoken)}{টোকেন (ṭoken)}]{\bn{টোকেন} (ṭoken) — a queue token}

\label{lesson:BN-C541-token}



\begin{quote}
Before the new one: say the Bengali for long-lasting, then the Bengali for reserved.
\end{quote}

\subsection*{You'll want to know: \bn{টোকেন}}
\textbf{\bn{টোকেন}} — \emph{ṭoken} — \textquotedblleft{}a queue token\textquotedblright{}.

\begin{quote}
\bn{টোকেন} \bn{না} \bn{নিলে} \bn{ব্যাংকে} \bn{কেউ} \bn{তোমাকে} \bn{ডাকবে} \bn{না।}
\end{quote}

\begin{quote}
\emph{Ṭoken nā nile bængke keu tomāke ḍākbe nā.}
\end{quote}

\begin{quote}
If you don't take a token, no one at the bank will call you.
\end{quote}

\begin{cousinweb}
\textbf{\bn{টোকেন}} (\emph{ṭoken}) is borrowed from English \emph{token}.
\end{cousinweb}

\begin{grammarlens}[title={What you've built}]
One more word for saying what happens if.
\end{grammarlens}

\subsection*{Your turn}
\begin{itemize}
\raggedright
  \item \emph{Say it:} \emph{ṭoken}
  \item \emph{Say it:} \emph{ṭoken}, once more
  \item \emph{Say it:} \emph{Ṭoken nā nile bængke keu tomāke ḍākbe nā.} — If you don't take a token, no one at the bank will call you.
  \item \emph{Recall:} say the Bengali for long-lasting, then the Bengali for reserved, then say \emph{ṭoken} again
\end{itemize}

\begin{tcolorbox}[breakable,colback=teal!4,colframe=teal!35!black,title={Before you move on}]
What does \bn{টোকেন} mean? (\textquotedblleft{}A queue token\textquotedblright{}.) Say it once more.
\end{tcolorbox}
\section[\texorpdfstring{\bn{পর্যায়} (pôrjāy)}{পর্যায় (pôrjāy)}]{\bn{পর্যায়} (pôrjāy) — a phase; a stage (of a process)}

\label{lesson:BN-C541-porjay}



\begin{quote}
Before the new one: say the Bengali for reserved, then the Bengali for a queue token.
\end{quote}

\subsection*{You'll want to know: \bn{পর্যায়}}
\textbf{\bn{পর্যায়}} — \emph{pôrjāy} — \textquotedblleft{}a phase; a stage (of a process)\textquotedblright{}.

\begin{quote}
\bn{প্রথম} \bn{পর্যায়ের} \bn{কাজ} \bn{শেষ} \bn{হলে} \bn{আমরা} \bn{ছুটি} \bn{নেব।}
\end{quote}

\begin{quote}
\emph{Prothom pôrjāyer kāj shesh hôle āmrā chhuṭi nebô.}
\end{quote}

\begin{quote}
When the first stage of the work is done, we'll take a holiday.
\end{quote}

\begin{cousinweb}
\textbf{\bn{পর্যায়}} (\emph{pôrjāy}) is a tatsama, Sanskrit \emph{paryāya}, which Monier-Williams gives as \textquotedblleft{}a turning round\textquotedblright{} and as \textquotedblleft{}succession, a turn\textquotedblright{}; the sense \textquotedblleft{}a stage, a phase\textquotedblright{} is the modern language's own.
\end{cousinweb}

\begin{grammarlens}[title={What you've built}]
One more word for saying what happens if.
\end{grammarlens}

\subsection*{Your turn}
\begin{itemize}
\raggedright
  \item \emph{Say it:} \emph{pôrjāy}
  \item \emph{Say it:} \emph{pôrjāy}, once more
  \item \emph{Say it:} \emph{Prothom pôrjāyer kāj shesh hôle āmrā chhuṭi nebô.} — When the first stage of the work is done, we'll take a holiday.
  \item \emph{Recall:} say the Bengali for reserved, then the Bengali for a queue token, then say \emph{pôrjāy} again
\end{itemize}

\begin{tcolorbox}[breakable,colback=teal!4,colframe=teal!35!black,title={Before you move on}]
What does \bn{পর্যায়} mean? (\textquotedblleft{}A phase; a stage (of a process)\textquotedblright{}.) Say it once more.
\end{tcolorbox}
